\documentclass{article}
\usepackage{import}


\title{Advanced Algorithms\\Homework 06}
\author{Felix Céard-Falkenberg\\7174020}

\usepackage{../homework} % See homework.sty %

\begin{document}

\section{Question 2}

\subsection{\textbf{Show that the \textit{Independent Set Problem} in planar graphs can be solved in $\mathcal{O}(6^k n)$ time using branching.}}

We know that each node of a planar graph has a degree of at most 5. 
% Therefore, we can test all possible combinations of nodes to include in the independent set 
Therefore we can test at each step, each of the 6 possible combinations of nodes to include in the independent set. 
Assuming that a branch has a depth of $k$, then we have found an independent set of size at least $k$.
Therefore, in the worst case, we have to test all possible combinations of nodes to include in the independent set, which is $6^k$. Lastly, because the graph might not be connected, we have to repeat this process for each of the connected components in the graph, which is then in the worst case $n$ times. Therefore, we have $\mathcal{O}(6^k n)$ time complexity.

\section{Question 3}
\subsection{\textbf{Show that the \textit{Independent Set Problem} in planar graphs has a linear sized kernel.}}

Given that every planar graph is at 4-colorable, we know that if $n \geq 4 \cdot k$, then we can always find an independent set of size at least $k$. In that case, we can return yes immediately.

Otherwise, we have $n < 4k$. In that case, we leave the instance unchanged. Hence the remaining instance has at most $4k-1$ vertices, which is linear in $k$.

% Thus, we obtain in polynomial time an equivalent instance $(G',k')$ such that either it is a trivial yes-instance, or it has size at most $4k-1 = \mathcal{O}(k)$. Therefore, the Independent Set Problem on planar graphs has a linear sized kernel.




\end{document}
